\section{Introduction}\label{sec:intro}
World models are designed to discern the current state of the environment and predict subsequent states based on executed actions. 
%Through simulation of prospective scenarios, world models facilitate a comprehensive understanding of environmental evolution under varying conditions.
This predictive ability of world models not only enables the assessment of decision-making consequences but also facilitates the generation of synthetic data, which serves as a crucial resource for training, testing, and simulation in autonomous systems. 
Consequently, world models have emerged as a focal point of research in the field of autonomous driving, garnering substantial attention from both academia and industry.

The synthetic data generated by world models can take various forms to represent future scenes.
%These include 2D videos that visually depict future environments~\cite{hu2023gaia, uniscene, wang2024drivedreamer, wang2024driving},
These include 2D videos~\cite{hu2023gaia, uniscene, wang2024drivedreamer, wang2024driving},
%3D lidar point clouds that offer a detailed spatial representation of the surroundings~\cite{agro2024uno, khurana2023point, uniscene, yang2024visual, zyrianov2024lidardm},
3D lidar point clouds~\cite{agro2024uno, khurana2023point, uniscene, yang2024visual, zyrianov2024lidardm},
%and 3D occupancy grids that provide a discretized representation of the space occupied by objects~\cite{occworld, occvar, dome,dfit-occworld,occ-llm, uniscene, renderworld, occllama}.
and 3D occupancy grids~\cite{occworld, occvar, dome,dfit-occworld,occ-llm, uniscene, renderworld, occllama}.
Irrespective of the representation format, the changes in the appearance of future scenes predicted by world models are predominantly governed by two key factors:
1) Ego-vehicle motion: The movement of the autonomous vehicle alters its viewing perspective, leading to dynamic changes in perceived spatial features (e.g., perspective shifts, occlusions).
2) Scene evolvement: Natural changes in the environment, such as agent interactions (e.g., pedestrian movements, vehicle collisions) and background updates (e.g., traffic light changes, weather variations).
%The first factor, ego-vehicle motion, pertains to the alterations in the visual and spatial information perceived by an autonomous vehicle as it traverses the environment.
%As the vehicle moves, the perspective from which it observes the world changes, leading to alternations in the model's representation.
%The second factor, scene evolvement, encompasses changes arising from modifications in the background environment and interactions among various agents, such as vehicles and pedestrians.
Notably, these two sources of change exhibit distinct characteristics. In a wide range of real-world scenarios, the background environment may remain relatively stable, with minimal changes over time (see \cref{fig:compare_video_occ}).
Instead, it is the motion of the ego-vehicle and the resulting shifts in the viewing perspective that contribute significantly to the dynamic changes in the world model's representation.

\begin{figure}
    \centering
    \includegraphics[width=0.9\textwidth]{fig/come_intro.pdf}
    \caption{\mtd{} with both scene-centric and ego-centric representation. Compared to ego-centric evolution, scene-centric prediction shows smaller gap in the context of temporal evolution. \mtd{} uses scene-centric prediction results as an important guidance to enhance occupancy world model.
    }
    \label{fig:compare_video_occ}
    % \vspace{-1em}
\end{figure}

State-of-the-art (SOTA) world models, however, rely on neural networks to implicitly learn the intertwined effects of these factors, often resulting in suboptimal spatial consistency. As a case study, the SOTA model DOME~\cite{dome} achieves an average 27.10 mIoU for 3-second occupancy prediction using ground-truth trajectories. In contrast, when predictions are formulated in a scene-centric coordinate system (where static backgrounds are decoupled from ego motion), a vanilla U-Net achieves 39.12 mIoU—a 44\% improvement. This disparity underscores the critical lack of explicit control over spatial consistency in existing methods, motivating the need for disentangled representation learning.

%However, a critical analysis of state-of-the-art world model methodologies reveals that they predominantly rely on neural networks to implicitly learn the intertwined representation changes associated with both factors simultaneously.
%In this study, we hypothesize that separating these two fundamental components can significantly mitigate the challenges inherent in world model tasks.
%For instance, the state-of-the-art DOME~\cite{dome} can only achieve an average 27.10 mIoU for the 3-seconds prediction task, with ground truth future trajectories.
%In constrast, if we move the prediction into a scene-centric coordinate system, the inherent world model difficulties can be decoupled as the static background environments no longer changes with ego-motion. Under such scene-centric settings, a vanilla U-Net can achieve a 39.12 mIoU, an astonishing 44\% improvement.
%This significant gap highlights the lack of explicit spatial consistency control in existing world models, which motivates us to introduce such control by decoupling environmental dynamics from ego-motion into scene-centric representations.



%%%%%%%%%%%%%%%%%%%
To address this, we propose \mtd{}, a three-stage framework that leverages scene-centric coordinates to separate ego motion from scene dynamics (see \cref{fig:method}):
1) Pose-conditioned generative stage: A diffusion-based model generates future occupancy maps by iteratively denoising latent vectors, using historical occupancy encodings and pose/BEV layouts as conditional inputs.
2) Fixed-view forecasting branch: By transforming past and future frames into a common coordinate system, this stage mitigates ego-motion effects, enabling non-interactive occupancy prediction without explicit scene flow estimation.
3) Global denoising diffusion stage: Cascaded Diffusion Transformer~\cite{DiT} blocks extract spatio-temporal features, with skip connections ensuring cross-layer consistency.
The Model Transfer Design (MTD) acts as the core mechanism, transferring knowledge from fixed-view forecasts to variable-view generation.
Structured as a trainable copy of the generative model's first half, it injects scene-centric condition features into the latter half via skip connections, enhancing generative realism and temporal consistency.

%\textbf{A bit verbose, revisit later}.
%Without loss of generality, we adopt 3D occupancy as the output representation for our world models.
%To verify our hypotheses, we propose a three-stage occupancy world model, as illustrated in \cref{fig:method}.
%1) In the first stage, the pose-conditioned generative model operates analogously to a diffusion-based image generation framework. Here, the historical segments of the noise latents are substituted with occupancy latents encoded from historical occupancy sequences.
%The model iteratively denoises the future components of the noise latents, which are subsequently decoded to generate future occupancy maps.
%Pose information and optional Bird's-Eye-View (BEV) layouts are incorporated as conditional inputs to the model blocks.
%2) The second stage features a fixed-view occupancy forecasting branch. Leveraging the dynamic characteristics of the majority of static and minor scene elements, this branch conducts non-interactive occupancy prediction by transforming past and future frames into the current frame's coordinate system.
%This approach eliminates the need for explicit scene flow decoupling and mitigates the impact of self-driving motion on predictions.
%3) In the final stage, the COME world model integrates a global denoising diffusion architecture composed of Diffusion Transformer~\cite{DiT} blocks.
%These cascaded DiT blocks sequentially extract spatio-temporal features, with positive layers connected to their reciprocal counterparts.
%The \mtd{} ControlNet serves as the pivotal mechanism for transferring knowledge from fixed-view forecasting to variable-view generation. Structurally, it is the trainable copy of the first half of the occupancy world model and conveys controls features to the second half of the world model with skip connections.


Extensive experiments on the nuScenes-Occ3D dataset validate \mtd{}'s efficacy across diverse settings.
\begin{inparaenum}[a)]
\item Input sources: ground-truth occupancy, camera-based predictions, and fusion-based occupancy inputs;
\item Prediction horizons: short-term (3s) and long-term (8s) forecasts;
\item Trajectory types: ground-truth poses and end-to-end planned trajectories.
\end{inparaenum}

In all configurations, \mtd{} outperforms SOTA baselines by significant margins, with scene-centric settings demonstrating the largest gains. Key contributions include:
\begin{inparaenum}[1)]
\item A divide-and-conquer strategy for occupancy world models, decomposing complex spatio-temporal prediction into ego-motion-agnostic and scene-dynamic components;
\item \mtd{} ControlNet, which leverages scene-centric forecasts as guidance to improve generative fidelity and spatial consistency;
\item Empirical validation of disentangled representation learning, establishing new state-of-the-art performance on a challenging autonomous driving benchmark.
%\item We propose a divide-and-conquer strategy for the occupancy world model, which makes it easier to understand the scene evolution with ego actions.
%\item We propose COME ControlNet, which first uses scene-centric forecasting to understand the evolution of the environment, and then injects it into the generative world model as a guidance to improve the realism and consistency of the generation.
%\item Extensive experiments on nuScenes-Occ3D benchmark under multiple settings show that the proposed \mtd{} significantly enhances the occupancy generation performance to the state-of-the-art.
\end{inparaenum}


%The effectiveness of our hypothesis is validated across multiple configurations of the nuScenes-Occ3D dataset. COME consistently outperform all prior arts with large margins.
% We conduct extensive evaluations of our model across multiple configurations on the nuScenes-Occ3D dataset.
%Specifically, we assess the model using 3D occupancy inputs derived from ground-truths, as well as those generated by camera-based and fusion-based 3D occupancy prediction methods.
%Our experiments consider both ground-truth poses and end-to-end planned trajectories, and span prediction horizons of 3 seconds and 8 seconds.
% Additionally, we benchmark our model against state-of-the-art baselines that utilize BEV layouts as optional inputs for simulation tasks.

% We also analyze the limitation of current metrics with visualization results and propose new evaluation protocols.
% The video world model uses a single image or historical video as a control quantity, and can freely generate future videos to understand the evolution of the world.

% The video world model uses a top-down thumbnail as a control quantity, which can control the generation of simulation synthetic data from multiple perspectives, and the maps and obstacles that meet the established requirements



%occupancy is one of its typical input/output representations, occupancy grids have ideal properties, such as precise geometric properties, rich semantic meanings, and easy to integrate with both learning-based and classical planners.

%the appearance changes of occupancy/camera images/lidar point clouds, are due to two major factors, one is ego-motion, one is scene changes due to background environment changes, together with agent interactions

%if we do the math/statiatics on major datasets, we noticed that background environment entail a very large portion of the voxels, decoupling the background environment changes from more important active agent foreground can greatly ease the job of the world model.
%by moving the predicting coordinate system to a scene-centric system, we noticed great/huge improvement for occ tasks

%however just do the scene-centric representation is not enough, occlusions and other problems, then we need to add controlnet to solve these problems

%by doing these things together, we conjueture that doing this in the scene centric paradigm can greatly improve the metrics. we want the community to

%potential problems: long horizon predicition at high speed, the basic assumption that most voxels stay put might no longer hold
%potential problems: current metric treat all voxels eaqually, we might need metrics that favor voxels/agents that interacts with the ego vehicle, eg, things behind the ego vehicle might not need to be considered